\xchapter{Five Letters by Cardinal Newman}

\textbf{[}to
John Rickards Mozley, son of Newman's sister, Jemima—\textbf{NR}.\textbf{]}

[from \emph{The
Contemporary Review}, 76, Sept. 1899. Also available in \emph{Letters
and Diaries of John Henry Newman}, vol. 27.]

THE
following letters of my uncle, the late Cardinal Newman, written in
the year 1875 (three years before he was made Cardinal), form
one—that is, the defensive—side of a correspondence which I had
with him in the course of that year. I had ventured to put to him
certain questions, to which he (as I felt sure would be the case) was
willing to reply. I had asked whether the real conduct of the visible
Church—\emph{i.e}., in his view, of the Church of Rome—had been in
accordance with that spirit of morality and goodness which should mark
a divine example and a divine teacher. I pointed to facts in the
history of the Church which appeared to me to be symptoms of a faulty
nature. I referred to the condition of the countries most obedient to
Rome—Spain under Philip II., France up to the first Revolution,
Italy up to the middle of the nineteenth century—as exhibiting a
tremendous total of misdoing, partly traceable directly to the
influence of the highest authorities of Rome, partly permitted by them
without protest or repudiation. How came it that the visible happiness
and harmony of the several countries of Europe should be almost in the
reverse proportion to the degree of their belief in the authority of
Rome? How came it that the members of an organization, to which the
divine promises were believed to have been entrusted, should not only
have committed such grave offenses in the past, but should be so
unwilling to confess them in the present, except as bare facts, and
without any sense of the disrepute thereby attaching to themselves,
and to the society they looked upon as divine?

The Cardinal's answers to the questions of which
the above is a summary will certainly be found extremely interesting.

With respect to my own share in the
correspondence I have but one regret; that, however, is a serious one:
it is that in relying to the Cardinal's last letter (that dated
December 3, 1875) I was overpowered by the magnitude of the subject,
and perhaps also by the personality of my opponent in argument, and
missed the true point. Hence, this was the end of continuous and
sustained argument between us; though, of the letters which I had from
him in later years, two certainly are of very high interest.

Let me, as far as is possible, repair my error by
indicating the central point in the Cardinal's letter of December 3,
1875, from which, were the opportunity ever to offer, the
argument might be resumed. It lies in the following sentence: ``The \emph{ethos}
of the Catholic Church is what it was of old time, and whatever or
whoever quarrels with Catholicism now, quarrels virtually, and would
have quarrelled if alive 1800 years ago, with the Christianity of
Apostles and Evangelists.''

The question, it will be seen, is this—and
truly it is an important one—whether the spirit of St. Peter and St.
Paul can be shown to differ, in any material respect, from the spirit
of the Church of Rome at the present day.

J. R. MOZLEY.

———————

\textbf{[}See also letter
of April 19, 1874, in Ward, vol. 2, pp. 572-574—\textbf{NR}\textbf{]}

THE ORATORY,

\emph{April} 1, 1875.

MY DEAR
JOHN,

You open a subject too large to be dealt with in one letter; but I
shall be able to get a certain way in it today.

I consider your letter to be addressed to me
personally, as if you said, ``I am perplexed and even curious to
understand how a man like you, who have had time and opportunities for
observation and thought, should be able to put up with a one-sided
view of the Church of Rome—nay, with an abstract view and a paper
representation of it, a mere conclusion, congruous or compulsory, from
premises dependent on certain first principles, such as 'there must be
a visible Church,' instead of going into the world of facts, and
seeing and judging of the Roman Church by what it is seen to have been
in history.''

My reply to your objections, then, shall take the
shape of accounting for my own insensibility to them as objections.
But anyhow, as I can only answer you in my own way and from my own
standpoint, the substance of what I shall say would be the same,
whether I argued with you directly or explained to you the arguments
which convince \emph{me}.

First, then, I grant that I do assume certain
first principles as the starting points from which my convictions
proceed, and I don't see who can arrive at any conviction without
making assumptions. I assume that there is a truth in religion, and
that it is attainable by us: that there is a God, to whom we can
approve ourselves and to whom we are responsible. On the other hand, I
find, in matter of fact and by experience, that there are great
difficulties in admitting this first principle; but still, they are
not such as to succeed in thrusting it out from its supremacy in my
mind. The most prominent difficulty of Theism is the existence of
evil: I can't overcome it; I am obliged to leave it alone, with the
confession that it is too much for me, and with an appeal to the
argumentum ab ignorantiâ, or, in other words, with the evasion or
excuse, not very satisfactory, that we have not the means here of
answering an objection, which nevertheless, if we knew more, we should
doubtless have the means of answering: that we can at least make
hypotheses to help the difficulty, and, though all those which
we can make be wrong, still they open a possibility and prospect of
other hypotheses as yet unknown, one of which may be the true
explanation.

When I come to Christianity I find this grand
difficulty untouched, yet fully recognised. This coincidence is to me
an argument in favour of Christianity, if Theism be true, as falling
under the argument from analogy. And, though Theism were not yet
proved true, still, from the fact of the coincidence, an argument in
some sort is to be drawn in favour of both systems, that is, supposing
the coincidence is independent of themselves—I mean, if Theists and.
Christians have not borrowed their recognition and non-explanation of
the fact of evil from each other.

Our Lord's death to destroy evil is as tremendous
and appalling a confession of the (its?) existence and of its power as
can be conceived.

From this central doctrine of the Gospel, the
Atonement, may be drawn two contrary conclusions. The first is that
from the moment of our Lord's death upon the cross all evil would be
annihilated; or secondly, that since He did not in His own Person
destroy it instantaneously, no wonder if He should take time in
destroying it in the world or in His Church. The former of those
conclusions is perhaps the more natural; but the interval of gloom and
sadness which overwhelmed His followers on His death, and still more
their history, as contained in the Acts of the Apostles, is sufficient
to show that it is not the right conclusion.

I confess, then, that it was natural, very
reasonable, to expect that an annihilation of sin and a millennium
period would commence with our Lord's Sacrifice; but, unless we
unravel our convictions and run back to belief in nothing, I must give
this thought up, and must admit, on the contrary, the pregnant
conclusion that evil will pass away from this world and from the
Church very slowly—nay (if we are to imagine that the moral system
advances after the analogy of the advance in the physical system of
the universe), so slowly that one or two generations or centuries
afford no available measure for calculating the rate of advance. I own
I should have fancied, \emph{a priori}, that the lamb and the lion
would lie down together from the date of the Crucifixion; that at
least that Elect Society which our Lord left behind Him would show
forth in its extension as a kingdom of righteousness from the first,
simple and absolute holiness extending with its extension, whereas, in
fact, the history of the Church contains in it the history of great
crimes.

I allow, then (and for argument's sake I allow
more than facts warrant), the existence of that flood of evil which
shocks you in the visible Church; but for me, if it touched my faith
mortally in the divinity of Catholicism, it would, by parity of
reason, touch my faith in the Being of a Personal God and Moral
Governor. The great question to me is not what evil is left in
the Church, but what good has energised in it and been practically
exercised in it, and has left its mark there for all posterity. The
Church has its sufficient work if it effects positive good, even
though it does not destroy evil except so far forth as it supplants it
for good.

Of its greatest and best achievements it cannot,
from the nature of the case, leave memorials, that ``hidden man of the
heart'' of which I spoke in that former letter to which you refer. It
is not necessarily seen in school teachers or in every specimen of a
secular priest, even though, did you know them, you might find that
your first impressions had been unjust to them. Nay, I have always
laid great stress on St. Paul's words, ``I endure all \emph{for the elects'
sake}''; they lead me to reflect that, even though there were no
high religious fruits of the Church's special sacraments generally,
ordinarily and \emph{primâ facie} visible to the world, that would
not necessarily be a refutation of its claim to come from God. The
Church would indeed, if it had no visible tokens at all, be a secret
society; but, since it is a light set on a hill, I grant it must have
visible tokens that it is divine, and, contrariwise to what you hold,
I think that it and its tokens are visible for the very reason that
God is invisible—viz., because they are to manifest Him. However,
though I grant that there must be visible tokens of sanctity in the
Church if the Church is to be considered divine, still, as the Spirit
bloweth as it listeth, so its manifestation in works is according to
no law and cannot be reckoned on.

As to the virtues of Catholics, I have lately
been reading the following words of Lord Russell, an impartial
witness, from his ``Essay on the Christian Religion'': ``There is among
Roman Catholics, in their relations to each other, a pure essence of
affection which does not appear in the moral writings of Greece and
Rome. The Roman Catholics, who have never practised or have
relinquished the vices of erring youth, are humble, loving,
compassionate, abounding in good works, kind to all classes of their
fellow creatures, ever ready to say, 'God be merciful to me a sinner,'
ready to give of their substance to the needy, ready to forgive others
their trespasses, and kneel in humble devotion to their Maker.'' He
speaks as if there were no middle class among us; but, if we were not
living in sin, we were almost saints.

But leaving the highest and truest outcome of the
Catholic Church and descending to history, certainly I would maintain
firmly, with most writers on the Evidences, that, as the Church has a
dark side, so (as you do not seem to admit) it has a light side also,
and that its good has been more potent and permanent and evidently
intrinsic to it than its evil. Here, of course, we have to rely on the
narrative of historians, if we have not made a study of original
documents ourselves. It would be a long business (assuming their
correctness), but an easy business too, to show how Christianity has
raised the moral standard, tone, and customs of human society;
and it must be recollected that for 1500 years Christianity and the
Catholic Church are in history identical. The care and elevation of
the lower classes, the championship of the weak against the powerful,
the abolition of slavery, hospitals, the redemption of captives,
education of children, agriculture, literature, the cultivation of the
virtues of piety, devotion, justice, charity, chastity, family
affection, are all historical monuments of the influence and teaching
of the Church. Turn to the non-Catholic historians, to Gibbon, Voigt,
Hurter, Guizot, Ranke, Waddington, Bowden, Milman, and you will find
that they agree in their praises, as well as in their accusations, of
the Catholic Church. Guizot says that Christianity would not have
weathered the barbarism of the Middle Age but for the Church. Milman
says almost or altogether the same. Neander sings the praises of the
monks. Hurter was converted by his historical researches. Ranke shows
how the Popes fought against the savageness of the Spanish
Inquisition. Bowden brings out visibly how the cause of Hildebrand was
the cause of religion and morals. If in the long line there be bad as
well as good Popes, do not forget that long succession, continuous and
thick, of holy and heroic men, all subjects of the Popes, and most of
them his direct instruments in the most noble and serviceable and most
various works, and some of them Popes themselves, such as Patrick,
Leo, Gregory, Augustine, Boniface, Columban, Alfred, Wulstan, Queen
Margaret of Scotland, Louis IX., Vincent Ferrer, Las Casas, Turibius,
Xavier, Vincent of Paul—all of whom, as multitudes besides, in their
day were the life of religion.

I have hardly begun my answer to your question,
yet I have written all this—but it is hard to be short on such a
subject. I shall stop here, and hope in a few days to come to closer
quarters with your main difficulty.

Yours affectionately,

JOHN H. NEWMAN.

THE ORATORY,

\emph{April} 4, 1875.

MY DEAR
JOHN,

Thank you for your letter of this morning, which leads me to say that
I did not use the word ``curious'' in a sense inconsistent with
earnestness in inquiry, though I cannot be sorry for an accident which
has been the occasion of your sending me so frank and \emph{ex animo}
an explanation.

I wish I could be shorter, but it is easier to
ask than to answer questions. In what I wrote to you the other day I
said that both good and bad were to be expected in the Catholic
Church, if it came from our Lord and His Apostles, whereas you
had ignored the good altogether, and had insisted there was in it an
actual tradition or abiding system of bad, forming a whole and giving
the Church a character; and worse, that, though it was so, Catholics
would not confess it and renounce it. Now I do confess that bad is in
the Church, but not that it springs from the Church's teaching or
system, but, as our Lord and His Apostles predicted it would be, in
the Church, but not of it. He says, ``It must \emph{needs be} that
scandals come;'' ``many are called, few are chosen;'' ``the kingdom of
heaven is like a net which gathereth of every kind.'' Good men and good
works, such as we find them in Church history, seem to me the
legitimate birth of Church teaching, whereas the deeds of the Spanish
Inquisition, if they are such as they are said to be, came from a
teaching altogether different from that which the Church professes.

It is on the Inquisition that you mainly dwell;
the question is whether such enormity of cruelty, as is commonly
ascribed to it, is to be considered the act of the Church. As to Dr.
Ward in the \emph{Dublin Review}, his point (I think) was not the
question of \emph{cruelty}, but whether persecution, such as in Spain,
was \emph{unjust}; and with the capital punishment prescribed in the
Mosaic law for idolatry, blasphemy, and witchcraft, and St. Paul's
transferring the power of the sword to Christian magistrates, it seems
difficult to call persecution (commonly so called) \emph{unjust}. I
suppose in like manner he would not deny, but condemn, the \emph{craft}
and \emph{cruelty}, and the wholesale character of St. Bartholomew's
Massacre; but still would argue in the abstract in defence of the
magistrate's bearing the sword, and of the Church's sanctioning its
use, in the aspect of \emph{justice}, as Moses, Joshua, and Samuel
might use it, against heretics, rebels, and cruel and crafty enemies.

I think such insane acts as St. Bartholomew's
Massacre were prompted by mortal fear. The French Court considered
(rightly or wrongly) that if they did not murder the Huguenots, the
Huguenots would murder them. Thus I explain Pope Gregory's hasty
approbation of so great a crime, without waiting to hear both sides.
After a period of luxury and sloth, the sudden outburst of the
Reformation frightened the Court of Rome out of its wits, and there
were those who thought the one thing needful was to put it down
anyhow, as the destruction, at least eventually, of all religion,
morality, and society. Perhaps they were right in this fear; and thus
they got mixed up with mere politicians, unscrupulous men, and became
in the eyes of posterity answerable for deeds which were not properly
theirs. I was reading the other day a defence of Pius V. against Lord
Acton, the point of which was that in no sense was it the Pope who
sanctioned the plot for assassinating Elizabeth, but, the Duke of
Alva. Yet who can deny, true as this may be, still that to readers of
history the Pope and the Duke are in one boat? Then, again, their
agents, or the sovereigns who sought their sanction for certain
courses or measures, went far beyond the intention of the Popes, who
nevertheless, from their political entanglements, could not resume the
powers that they had once given over to them. A large society, such as
the Church, is necessarily a political power, and to touch politics is
to touch pitch. A private Catholic is not answerable for the Pope's
political errors, any more than the shareholder in a railway in 1875
is answerable for the railway's accidents in 1860, nay, or in 1875.

You say that at least the Popes ought publicly to
confess, when it is proved they have gone wrong. Does Queen Victoria
confess the sins of George IV.? Do principals feel it generous to
abandon their subordinates, or loyal children acquiesce in attacks on
their parents? As to controversialists, they are pleaders at a bar,
and are afraid to make admissions lest these should be turned against
them. To speak out is in the long run the wisest, the most expedient,
the most noble policy; seldom the possible, or the natural. Why are
private memoirs kept back from publication for thirty or sixty years?
No party can be kept together if there is no reticence. But in fact,
except among controversialists, there is no want of candour and
frankness among us; witness the fact that Protestant attacks on us
generally are drawn from the admissions of Catholics. Baronius,
writing under the Pope's eye, speaks in the strongest terms of the
evil state of the Popedom in the dark age; Rinaldus, his continuator,
speaks against Alexander VI.; St. Bernard, St. Thomas, and many others
speak against the conduct of the Roman See in their own times. So do
Pope Adrian VI., Paul IV., \&c. So do holy women in their writings,
such as St. Bridget.

As to the state of Catholic Europe during these
last three centuries, I begin by allowing or urging that the Church
has sustained a severe loss, as well as the English and German
nationalities themselves, by their elimination from it; not the least
of the evil being that in consequence the Latin element, which is in
the ascendant, does not, cannot know, how great the loss is. This is
an evil which the present disestablishment everywhere going on may at
length correct. Influential portions of the Latin races may fall off;
and if Popes are chosen from other nationalities, other ideas will
circulate among us and gradually gain influence.

As to the unbelief of France, Italy, and Spain,
allowing it to the extent facts warrant, still I had fancied that \emph{England},
the most fiercely Protestant country of Europe, had begun the
tradition of infidelity in Europe in its school of Deists in the
seventeenth and eighteenth centuries; and that \emph{Germany}, the
native soil of the Reformation, was now the normal seat of
intellectual irreligion. Is it not something the case of the pot and
the kettle?

Next, as to the bad government in the Papal
States, I allow, or rather argue, that an ecclesiastical
world-wide sovereign has neither time nor thought to bestow on secular
matters, and that such matters go to rack and ruin, and cause great
scandal in public opinion, as surely as would happen if I undertook to
be Chancellor of the Exchequer. The shortness of the reigns of the
Popes, an advantage ecclesiastically, and their political troubles,
increase this evil. Another thing—till of late there was no science
of government, and the Papal administration was not worse than its
neighbours; but now we have a dozen sciences, political, economical,
sanitary, social, agricultural, municipal, and the like, all tending
to the tranquillity and prosperity of States, which secular
Governments can carry out and profit by, but which ecclesiastics and
theologians have no head for.

Further, all States have their course, their
beginning and their end. It is not wonderful that those which were
great three centuries ago should be waning or dying out now; while
England, which then was barbarous compared with the Continent, and
much more Prussia, Russia, and the United States, should be in the
ascendant. There is nothing to show what the state of England will be
two centuries hence. Its want of coal \emph{may} be its ruin; or,
before that want is felt, Protestantism, which has made it great, may,
by running into democracy, make it small again. At present the
Catholic Church is encumbered by its connection with moribund nations,
and, so far, Keble's application of the ``Mortua quinetiam,'' \&c.,
may be transferred to \emph{it}. Catholics are certainly taken at
great disadvantage now; but, as a loyal servant of Alfred or Bruce,
knowing the greatness of his master's soul and the splendour of his
gifts, might have no temptation whatever to mistrust his ultimate
success, in spite of temporary disaster, so we feel about the defects
and humiliations of the Papacy.

You see all along I have kept to my purpose of
describing my own view of the difficulties of Catholicity on which you
fasten, instead of attempting to deal with them controversially. The
temporal prosperity, success, talent, renown of the Papacy did not
make me a Catholic, and its errors and misfortunes have no power to
unsettle me. Its utter disestablishment may only make it stronger and
purer, removing the very evils which are the cause of its being
disestablished.

I was rejoiced to be told by you that you
recognised the truth of the power of prayer. Nothing else will clear
our religious difficulties.

Yours affectionately,

JOHN H. NEWMAN.

THE ORATORY,

\emph{April} 21, 1875.

MY DEAR
JOHN,

There is nothing ungenerous, as you fear, in your new questions; and,
if you had asked them distinctly before, I should have answered them
to the best of my power.

You now ask me whether I agree or disagree with
your judgment ``that the Church of Rome, as a society, has sometimes
done, more often sanctioned, actions, which were wrong and injurious
to mankind.'' I find no difficulty in answering you. I should say that
the Church has two sides, a human and a divine, and that everything
that is human is liable to error. Whether, so considered, it has in
matter of fact erred must be determined by history, and, for the very
reason that it is human as well as divine, I am disposed to believe it
has, even before the fact has been proved to me from history. At the
same time I must add that I do not quite acquiesce in the wording of
your question. It sounds awkward to ask, \emph{e.g}., ``Has the Kingdom
of England done or sanctioned wrong?'' It would be more natural to say,
``Has the nation done wrong, or the sovereign, or the legislature done
wrong, or all of these together''? I have no difficulty in supposing
that Popes have erred, or Councils have erred, or populations have
erred, in human aspects, because, as St. Paul says, ``We have this
treasure in earthly vessels,'' speaking of the Apostles themselves. No
one is impeccable, and no collection of men.

I grant that the Church's teaching, which in its
formal exhibitions is divine, has been at times perverted by its
officials, representatives, subjects, who are human. I grant that it
has not done so much good as it might have done. I grant that in its
action, which is human, it is a fair mark for criticism or blame. But
what I maintain is, that it has done an incalculable amount of good,
that it has done good of a special kind, such as no other historical
polity or teaching or worship has done, and that that good has come
from its professed principles, and that its shortcomings and omissions
have come from a neglect or an interruption of its principles.

The question that remains is, Has that which
claims to be divine in the Church sanctioned that which is human and
faulty in it? I maintain, No: and, in alleged cases brought in proof
of the affirmative, I should contend either that its sanction of the
act in question had no claim to be considered divine, or that the act
itself was not faulty. Thus St. Paul says, ``I wist not that he was
high priest, for it is written,'' \&c., and some commentators say
that he \emph{was} ignorant—that is, his act did not proceed from
the divine inspiration with which he was gifted; others that his act
was not wrong, for the man whom he reviled, in fact, was not high
priest.

However, I cannot simply grant to you, as you
assume, that mere omission to pronounce upon a faulty act is
necessarily itself a fault. Things are so constituted in this world,
that the power of doing good has a maximum. The Church, viewed as a
political body, has always been in advance of the age; up to 1600 most
men would grant this; but, as the Jews were allowed divorce as
practically a necessity in order to avoid worse evils, so it has not
always been possible for the Church to do upon the spot that which was
abstractedly best, as Elisha shirked the question of Naaman about
bowing in the house of Rimmon. Nor am I disposed to deny that, as time
goes on, the authoritative view of moral and religious truth becomes
clearer, wider, and more exact.

I do not know how I can answer your question more
closely than in what I have now said, and as, I think, I did answer it
in my former letter.

As to the last three centuries, the Church's
great battle has been against the various forms of error to which
Protestantism has opened the door. The work of the Church has on every
side been met and thwarted by the opposition of rival religions. In
India the work, begun by St. Francis Xavier, has been brought to a
stand by the variety and discordance of Christian sects. Still, if it
is a great work to preserve Christianity in the world, this I think
the Church has done and is doing: and at this moment Christianity
would be dying out in all its varieties were the Catholic Church to be
suppressed.

I hope I need not say I shall always feel a
pleasure and interest in hearing whatever you are moved to tell me
about yourself—pray do, for I am always

Yours affectionately,

JOHN H. NEWMAN.

P.S.—I hope you are out of anxiety by this time
about your little boy.

THE ORATORY,

\emph{May} 16, 1875.

MY DEAR
JOHN,

I am very glad to have so long a letter from you, but you must let me
wait, and be patient with me, as to my answering it, for I have
received a very heavy blow in the sudden and alarming illness of the
greatest friend I have—an illness, the issue of which will take some
time to show itself, and which has almost turned my head \footnote{Ambrose St. John—\textbf{NR}.}.

Thank E—— for wishing to send me, and you for
sending, her love—and tell her that I am very grateful to her, and
send her a double measure in return—one in reciprocity, and one from
gratitude.

Yours affectionately,

JOHN H. NEWMAN.

\emph{Dec}. 3, 1875.

MY DEAR
JOHN,

Your letter puts me into a great difficulty. It is my heart's desire
to bring you nearer to me in opinion, and so to explain my own
religious views as to excite in you interest and sympathy for us; to
reduce difficulties, and to inspire hope that Catholics and
Protestants are not so far apart from each other as is commonly said;
in a word, to throw myself into the sentiment which has led you to
write, and to co-operate with it. But I cannot feel you have gone to
the bottom of the matter, and it would not consist with that truth and
frankness due to all men, and especially to one with whom I am so
united in affection as yourself, not to say so.

I agree with you, then, but I go far beyond you
in holding, that the difference between Catholics and Protestants is
an ethical one; for I think that in pure Catholics and pure
Protestants (I mean, by so speaking, that most Protestants are tinged
with Catholicity, and most Catholics with Protestantism) this
difference is radical and immutable, as the natures of an eagle and a
horse are, except logically, two things, not one. Opposition to
physical science or to social and political progress, on the part of
Catholics, is only an accidental and clumsy form in which this vital
antagonism energises—a form, to which in its popular dress and
shape, my own reason does not respond. I mean, I as little accept the
associations and inferences, in which modern science and politics
present themselves to the mass of Catholics, as I do those contrary
ones, with which the new philosophy is coloured (I should rather say,
stained) by great Professors at Belfast and elsewhere.

Dealing with facts, not with imaginations,
prejudices, prepossessions, and party watchwords, I consider it
historically undeniable—

1. First, that in the time of the early Roman
Empire, when Christianity arose, it arose with a certain definite
ethical system, which it proclaimed to be all-important, all-necessary
for the present and future welfare of the human race, and of every
individual member of it, and which is simply ascertainable now and
unmistakable.

Next, I have a clear perception, clearer and
clearer as my own experience of existing religions increases, and such
as every one will share with me, who carefully examines the matter,
that this ethical system ([\emph{ethos}] we used to call it at Oxford
as realised in individuals) is the living principle also of present
Catholicism, and not of any form of Protestantism whatever—living,
both as to its essential life, and also as being its vigorous motive
power; both because without it Catholicism would soon go out, and
because through it Catholicism makes itself manifest, and is
recognised. Outward circumstances or conditions of its presence may
change or not; the Pope may be a subject one day, a sovereign another;
\emph{primus inter pares} in early times, the \emph{episcopus
episcoporum} now; there might be no devotions to the Blessed Virgin
formerly, they may be superabundant of late; the Holy Eucharist might
be a bare commemoration in the first century, and is a sacrifice in
the nineteenth (of course I have my own definite and precise
convictions of these points, but they are nothing to the purpose here,
when I want to confine myself to patent facts which no one ought to
dispute); but I say, even supposing there have been changes in
doctrine and polity, still the \emph{ethos} of the Catholic Church is
what it was of old time, and whatever and whoever quarrels with
Catholicism now, quarrels virtually, and would have quarrelled, if
alive, 1800 years ago, with the Christianity of Apostles and
Evangelists.

2. When we go on to inquire what is the ethical
character, whether in Catholicity now or in Christianity in its first
age, the first point to observe is that it is on all hands
acknowledged to be of a character in utter variance with the ethical
character of human society at large as we find it at all times. This
fact is recognised, I say, by both sides, by the world and by the
Church. As to the former of the two, its recognition of this
antagonism is distinct and universal. As regards \emph{Catholicism},
it is the great fact of this very day, as seen in England, France,
Germany, Italy, and Spain. On the other hand, we know that in the \emph{Apostolic}
Age Christians were called the ``\emph{hostes humani generis}'' (as the \emph{Quarterly}
called \emph{Catholics} within this two years), and warred against
them accordingly.

This antagonism is quite as decidedly
acknowledged on the side of the Church, which calls scciety in
reprobation ``the world,'' and places ``the world'' in the number of
its three enemies, with the flesh and the devil, and this in her
elementary catechisms. In the first centuries her badge and boast was
martyrdom; in the fourth, as soon as she was established, her war-cry
was, ``\emph{Athanasius contra mundum}''; at a later time her protests
took the shape of the Papal theocracy and the \emph{dictatus Hildebrandi}.
In the recent centuries her opposition to the world is symbolised in
the history of the Jesuits. Speaking, then, according to that aspect
of history which is presented to the eyes of Europeans, I say the
Catholic Church is emphatically and singularly, in her relation to
human philosophy and statesmanship, as was the Apostolic Church, ``the
Church militant here on earth.''

3. And, what is a remarkable feature in her \emph{ethos}
now and at all times, she wars against the world from love of it.
What, indeed, is more characteristic of what is called Romanism now
than its combined purpose of opposing yet of proselytisiug the
world?—a combination expressed in our liturgical books by the two
senses of the word ``\emph{conterere},'' that of grinding down and of
bringing to contrition. How strikingly, on the other hand, does this
double purpose come out in the Apostles' writings? We have three
primitive documents, each quite distinct in character from the
other two, differing in accidents and externals, but all intimately
agreeing in substantial teaching, so that we are quite sure of the
genius and spirit of Christian ethics from the first: I mean, (1) the
Synoptical Gospels, (2) St. Paul's Epistles, (3) St. John's Gospel,
Epistles, and Apocalypse. Now, the first of these says, ``Ye shall be
hated of all men for my name's sake. The disciple is not above his
Master. Fear them not. I came not to send peace on earth, but a sword.''
``I pray not for the world,'' says the third, ``the world hath hated them
because they are not of the world. Love not the world, neither the
things that are in the world. The world lieth in wickedness.'' And the
second, ``In time past ye walked according to the course of this world,
and were by nature the children of wrath even as the rest.'' And yet, ``Preach
the gospel to every creature,'' says the first; ``God so loved the
world, that,'' \&c., says the third; and ``He will have all men to be
saved,'' says the second. After avowals such as these in our primary
authorities, it will be a hard job to discover any Irenicon between
Catholicity and the moral teaching of this day.

4. This will be still clearer as we examine the
details of our ethics, as developed from our fundamental principles.
The direct and prime aim of the Church is the worship of the Unseen
God; the sole object, as I may say, of the social and political world
everywhere, is to make the most of this life. I do not think this
antithesis an exaggeration when we look at the action of both on a
large scale and in their grand outlines. In this age especially, not
only are Catholics confessedly behindhand in political, social,
physical, and economical science (more than they need be), but it is
the great reproach urged against them by men of the world that so it
is. And such a state of things is but the outcome of apostolic
teaching. It was said in the beginning, ``Take no thought for the
morrow. Woe unto those that are rich. Blessed be the poor; to the poor
the gospel is preached. Thou hast hid these things from the wise and
prudent. Not many wise men, not many mighty, not many noble are
called. Many are called, few are chosen. Take up your cross and follow
me. No man can have two masters; he who loveth father or mother more
than me is not worthy of me. We walk by faith, not by sight; by faith
ye are saved. This is the victory that overcometh the world, our
faith. Without holiness no man can see the Lord. Our God is a
consuming fire.'' This is a very different ethical system from that
whether of Bentham or of Paley.

5. I am far from saying that it was not from the
first intended that the strict and stern ethics of Christianity should
be, as it was in fact, elastic enough to receive into itself secular
objects and thereby secular men, and secular works and institutions,
as secondary and subordinate to the \emph{magisterium} of
religion—and I am far indeed from thinking that the teaching
and action of the world are unmixed evil in their first elements
(society, government, law, and intellectual truth being all from God),
and far from ignoring the actual goodness and excellence of individual
Protestants, which comes from the same God as the Church's holiness;
but I mean that, as you might contemplate the long history of England
or France, and recognise a vast difference between the two peoples in
ethical character and national life, and consequent fortunes, so, and
much more, you can no more make the Catholic and Protestant \emph{ethos}
one, than you can mix oil and vinegar. Catholics have a moral life of
their own, as the early Christians had, and the same life as
they—our doctrines and practices come of it; we are and always shall
be militant against the world and its spirit, whether the world be
considered within the Church's pale or external to it.

6. Coming back to your letter, I should not
wonder if you think I have mistaken its drift, and have been beating
the air. I do not think I have, though I have thought it best to fall
back upon the previous question.

I have meant to say that, though our opposition
to science, \&c., ceased ever so much, we should not thereby be
more acceptable in our teaching to the public opinion of the day.

Ever yours affectionately,

JOHN H. NEWMAN.

P.S.—Thank you for what you tell me of your new
abode. On reading this over I find it more difficult to follow its
course of thought than I could have wished.
